\chapter{Discussion}
\label{ch:discussionmain}

\section{What the architecture bought, and what it did not}

The degradation architecture kept the system running and its failures
visible: with eight of fourteen channels failing, the right arm retains
full position sensing and the left retains direction, and the incoherent
verdict exists because a failing potentiometer does not go quiet -- range
alone calls a wide-spanning nonsense signal healthy. What it did not buy is
capability: freezing the left arm's bend joints collapses its commanded set
from a shell to a surface, and no health scheme restores a measurement
never taken. The reduced-degree-of-freedom premise succeeded exactly where
a single sensor observes the quantity directly (elevation from gravity,
reach from a forward-kinematic magnitude) and failed where it does not: the
spherical model assumes azimuth and elevation are independent, and they are
not, because the first joint's axis does not tilt as the arm bends -- a
design rule for a successor, not a defect to patch with better calibration.

\section{The workspace result, and what it costs}

The disjoint-workspace finding is severe because it is geometric: no
planner or orientation policy can produce a point neither arm can reach, so
inter-arm handover was excluded from the task set before anything else
excluded it. The cause is not the mount's rotation but the arms'
asymmetric home joint angles, ground truth because the sim-to-hardware
bridge replays them onto the real arm; fixing it needed a mechanical change,
costed, applied and re-verified as one deliberate pass because every
clearance and task coordinate in the report depends on it.

\section{The safety layer is thorough about distance and silent about time}

The clearance floor is a \emph{distance}, and the wearer can move. Combining
the measured wearer-tracking latency with literature limb-speed ranges gives
\cref{tab:reaction-compact}: in most modelled limb motions the limb travels
further than the entire \SI{150}{\milli\metre} margin before the perception
chain is even aware anything changed -- not an implementation defect, but
the consequence of a fixed margin meeting a perception chain of roughly
half a second and a person who can move a metre in that time. It is this
report's principal open safety question.

\begin{table}[htbp]
  \centering
  \small
  \caption[The reaction budget, condensed]{Limb travel during the measured
  wearer-tracking latency against the \SI{150}{\milli\metre} clearance floor;
  full table and six limb-motion classes in \cref{ch:motion}.}
  \label{tab:reaction-compact}
  \begin{tabular}{@{}lll@{}}
    \toprule
    Latency & Deliberate reach & Quick reach or startle \\
    \midrule
    \SI{62}{\milli\second} (median) & \SI{62}{\percent} of the floor & exceeds the floor \\
    \SI{562}{\milli\second} (worst) & exceeds the floor & exceeds it several times over \\
    \bottomrule
  \end{tabular}
\end{table}

\section{Validity of the measurements}
\label{sec:validity}

Twenty-six times in this project a measurement was wrong and the system was
not (\cref{ch:instrument}), a high enough base rate to condition how every
other number here should be read. Fifteen of the twenty-six were caught by
an independently known quantity contradicting the reading rather than by
inspection -- the argument for building a cross-check into every
instrument rather than trusting it harder. Every result above except the
pilot and channel measurements comes from simulation with mock hardware,
whose force channels are identically zero, so any force-derived quantity is
structurally undefined; and every latency reported is control-path latency
between software stages, excluding the operator, display and physical arm,
so it is a floor.

\section{What the pilot shows, and its limits}
\label{sec:pilot-limits}

The pilot is the first evidence from real operators on this platform, and it
points the same direction as the shared-autonomy literature developed for
robots not attached to anybody~\cite{dragan2013policyblending,javdani2015hindsight}:
shorter completion times and fewer re-grips under assistance, with no
measurable cost in responsiveness. It should not be read as more than that.
Hand-travel distance is \emph{higher} under shared autonomy, and unlike
completion time this holds in the paired comparison too: six of seven
paired cases travel further under assistance, not less. This pilot cannot
say why -- plausibly the arbiter nudges a directed hand along a route that
is not the shortest the operator would have taken alone, making "shorter
but less direct" a real property of it rather than a contradiction, but
seven paired trials cannot separate that from task-order or fatigue. Five
operators with uneven task coverage is a pilot-scale sample, not a powered
comparison, and the pooled figures mix three tasks (plus P5's unlabelled
sessions) together; the paired comparison isolates the condition, and rests
on seven cases. Session duration includes setup and hesitation, which is
why hand travel and re-grips -- drawn from the position trace, not a
stopwatch -- are read as the more reliable measures. Most importantly, the
pilot has no separate wearer: it measures what the operator experienced,
not what carrying \SI{17}{\kilo\gram} of assisted motion costs a second
person, the comparison \cref{app:study} remains the target for.

\section{Against the planning report}

The planning report specified physical springs at the master and virtual
impedance at the slave. Neither survives contact with the measured system:
the arm's cyclic interface is unusable over the network path, so commands
go through a high-level interface at \SIrange{18.4}{18.7}{\hertz}, and an
impedance law at that rate is not one. The compliance sought would need to
come from the mechanism -- series elastic actuation at the
joint~\cite{toubar2022sea,lee2019sea} or a current-based formulation
avoiding a fast torque loop~\cite{dewolde2024impedance} -- both hardware
changes. The safety architecture, degradation architecture, workspace
characterisation and pilot were not in the plan; each followed from a
measurement the plan did not anticipate.
